\cotefile{157-5}
\scanpage{181}
% manuscript position: top
\mspage{97}
% manuscript position: 1–3
on pourra aussi bien traiter le cas général d'une \uline{$\lim$} de type $J$ quelconque ($J$ \unsure{petit|petite}). Pour ceci, soit
% manuscript position: 4–6
\[
\begin{aligned}
 I=I(J)={}&\text{ens des sous-catégories}\\
 &J'\text{ de }J\text{ telles que}\\
 &\operatorname{card}\operatorname{Fl}(J')\leq\pi,
\end{aligned}
\]
% manuscript position: 7–15
ordonnons $I$ par inclusion. Il est évident que $I\neq\emptyset$ (car $\emptyset\in I$), et \unsure{$=$} que \uline{$I$ est grand devant $\pi$}, car si $(J'_\alpha)_{\alpha\in A}$ avec $\operatorname{card}(A)\leq\pi$ est une famille d'\unsure{élmts|élémts} de $I$, alors la \inserted{sous-}catégorie \inserted{de $J$} engendrée par les $J'_\alpha$ est encore dans $I$ (vérification immédiate), et majore les $J'_\alpha$.
% manuscript position: 15–16
\deleted{On a \unsure{donc} pour $\varphi:J\longrightarrow\operatorname{Ind}'_\pi(C)$ donné, \unsure{considérons}, le foncteur}
% manuscript position: 17
\[
\deleted{I\longrightarrow\operatorname{Ind}'_\pi C\overset{\mathrm{déf}}{=}M}
\]

% manuscript position: 18–21
En fait, la considération de $I(J)$ associé à $J$ va nous permettre de prouver ceci, qui généralise 4.4.13 c) \unsure{rentre|contenu|rentrant} dans le b) :
% manuscript position: 22-margin
% slanted
\margin{\unsure{Comparer|Comp.} avec \unsure{énoncé} \unsure{4.10.3|5.10.3} page 270}

% manuscript position: 22–23
\uline{Corollaire 4.4.14} \quad Soit $M$ une catégorie, $\pi$ un cardinal infini tel que a) $M$ est
% manuscript position: foot
% archive sheet number: 180
